\chapter{Zwei Arten zu wachsen: Struktur oder Wert}
\label{ch:erweiterungen}

% Register: D5, S8. Quellen: Spez. Abschnitt 5 und 6.

\section{Der Moment der Gabelung}

Nehmen wir die Sieben in ihrem Zuhause $H_4$: dort schreibt sie sich
$0012$. Jetzt geben wir ihr eine Stelle mehr. Was soll das heißen?

Die erste Antwort ist die naive: Wir stellen eine Null davor.
\[
0012 \;\longmapsto\; 00012 .
\]
Die Ziffern bleiben stehen -- aber in $H_5$ wiegen alle Stellen
außer der letzten anders, und der Wert rutscht von $7$ auf $8$. Die \emph{Gestalt}
wurde bewahrt, der Wert nicht.

Die zweite Antwort ist die vorsichtige: Wir wollen, dass die Zahl
dieselbe bleibt, und codieren die $7$ in $H_5$ neu:
\[
0012 \;\longmapsto\; 00011 .
\]
Der \emph{Wert} wurde bewahrt, die Gestalt nicht.

Im Dezimalsystem sind beide Antworten dieselbe Operation -- deshalb
ist uns nie aufgefallen, dass es zwei Fragen sind. Im Horiraum
trennen sie sich, und wir geben ihnen Namen: Die strukturtreue --
gleichwertig sagen wir: gestalttreue -- Erweiterung $Z$ stellt die
Null voran, der werttreue Lift $L$ codiert um.

\begin{definition}[Register D5]
$Z_n(a_1,\ldots,a_n)=(0,a_1,\ldots,a_n)$ und
$L_n(a)=E_{n+1}(V_n(a))$, wobei $E_{n+1}$ die Codierung in
$\Hn{n+1}$ ist. Es gilt $V_{n+1}(L_n(a))=V_n(a)$, während $Z_n$ im
Allgemeinen den Wert verändert.
\end{definition}

\section{Zwei Treppen, zwei Himmel}

Beide Operationen verbinden die endlichen Horizonte zu einer
unendlichen Treppe -- aber es sind \emph{verschiedene} Treppen, und
sie führen an verschiedene Orte. Die Frage ist jeweils: Welche
Elemente aufeinanderfolgender Horizonte gelten als „dasselbe“?

Unter $L$ gelten zwei Darstellungen als dasselbe, wenn sie denselben
Wert haben: $013$ in $H_3$, $0012$ in $H_4$, $00011$ in $H_5$ sind
ein einziges Objekt -- die Sieben. Unter $Z$ gelten sie als
dasselbe, wenn sie dieselbe Gestalt haben: $12$, $012$, $0012$ sind
ein einziges Objekt -- das Muster „eins-zwei“, dessen Wert von
Stufe zu Stufe wandert.

\begin{satz}[Register S8]\label{satz:grenzen}
Der Grenzraum der Treppe unter $L$ ist $\Nn$, die Menge der
gewöhnlichen natürlichen Zahlen. Der Grenzraum unter $Z$ ist der
Raum aller Gestalten -- formal: der Halbring der Polynome mit
nichtnegativen ganzzahligen Koeffizienten in der Basis der
fallenden Fakultäten, den das nächste Kapitel einführt.
\end{satz}

\begin{proof}[Beweisidee]
Unter $L$ werden genau Darstellungen desselben Wertes identifiziert
(die Codierung ist injektiv), und jeder Wert kommt vor -- die
Klassen \emph{sind} die natürlichen Zahlen. Unter $Z$ werden genau
Ziffernfolgen identifiziert, die sich um führende Nullen
unterscheiden, und jede Gestalt kommt ab ihrem kleinsten zulässigen
Horizont vor -- die Klassen sind die Gestalten. Der vollständige
Beweis über gerichtete Systeme steht in Anhang~\ref{ch:anhang-beweise}.
\end{proof}

\begin{beispiel}
Die beiden Treppen, nebeneinander. Werttreu bleibt die Sieben sie
selbst und wechselt die Gestalt; gestalttreu bleibt das Muster
„eins-zwei“ und wechselt den Wert:
\begin{center}
\begin{tabular}{l|llll}
Welt & $H_2$ & $H_3$ & $H_4$ & $H_5$\\
\hline
werttreu (immer 7) & --- & $013$ & $0012$ & $00011$\\
gestalttreu (immer „12“) & $12=5$ & $012=6$ & $0012=7$ & $00012=8$\\
\end{tabular}
\end{center}
Man beachte die Kreuzung: $0012$ kommt in beiden Zeilen vor -- einmal
als zweite Wohnung der Sieben, einmal als drittes Leben des Musters.
Dasselbe Element, zwei Lesarten.
\end{beispiel}

Das ist der erste Satz dieses Buchs, den man philosophisch lesen
darf: \emph{Dieselben endlichen Räume erzeugen verschiedene
Unendlichkeiten, je nachdem, welchen Konsistenzbegriff man zwischen
den Stufen wählt.} Die vertraute Zahlengerade ist eine dieser
Unendlichkeiten -- die werttreue. Die andere, die gestalttreue, ist
mathematisch genauso legitim, nur hat sie im gewöhnlichen
Zahlensystem nie jemand von der ersten unterscheiden können, weil
dort beide Treppen zusammenfallen.

\section{Die dritte Treppe}

Lange dachten wir, es gäbe genau zwei Arten zu wachsen. Dann las ein
externer Mathematiker eine Kurzbeschreibung dieses Systems -- und
verankerte die Basenleiter versehentlich am anderen Ende. Aus seinem
„Missverständnis“ fiel eine dritte Treppe heraus.

Man kann einer Zahl nämlich auch \emph{rechts} eine Null anhängen.
Dabei behalten alle Ziffern ihre Position von links gezählt, und es
gilt eine verblüffend saubere Regel: Der Wert vervielfacht sich
nicht unkontrolliert, sondern skaliert \emph{exakt} mit dem Faktor $n+2$ --
aus $12$ in $H_2$ (Wert $5$) wird $120$ in $H_3$ (Wert $20$).
Invariant bleibt dabei weder der Wert noch die Gestalt, sondern eine
dritte Größe: der \emph{Fakultätsbruch}
$F=\sum_i a_i/(i+1)!$, denn es gilt stets
\[
V_n(a)=(n+1)!\cdot F(a).
\]
Jede Hori-Zahl ist also auch ein Bruch zwischen null und eins, mal
der Größe ihrer Welt -- und dieses Bruchsystem ist ein alter
Bekannter: Es ist Cantors Fakultätsbasis für Brüche, gut
anderthalb Jahrhunderte alt.

Damit hat der Horiraum drei Treppen, drei Invarianten, drei Himmel:

\begin{center}
\begin{tabular}{l|lll}
Erweiterung & invariant & Wert & Grenzraum\\
\hline
werttreu ($L$) & Wert & $\times1$ & natürliche Zahlen\\
gestalttreu ($Z$, links) & Polynom & ungleichmäßig & Strukturraum\\
bruchtreu ($R$, rechts) & Bruch $F$ & $\times(n+2)$ & Cantor-Brüche\\
\end{tabular}
\end{center}

Und der Einwand „das ist doch seit Cantor bekannt“ bekommt damit
eine präzise Antwort: Er trifft die dritte Treppe vollständig -- und
die anderen beiden samt allem, was dieses Buch über ihr Zusammenspiel
beweist, gar nicht.

\section{Der Eigenhorizont: wo die Gabelung heilt}

Für jede einzelne Zahl gibt es einen Punkt, an dem der Unterschied
zwischen $Z$ und $L$ verschwindet. Steht der gesamte Wert in der
letzten Ziffer -- die Zahl $n$ als $0\cdots0n$; wir nennen diese
Gestalt ihre \emph{Terminaldarstellung}, und sie ist ab dem
Horizont $H_n$ möglich --, dann ist die Nullerweiterung wertneutral:
Gestalt und Wert wachsen ab dort gemeinsam. Wir nennen $H_n$ den
\emph{Eigenhorizont} von $n$. Unterhalb davon wird die Zahl bei
jedem Umzug umgebaut; ab dort wächst sie nur noch still nach links.
Man kann exakt ausrechnen, wann eine führende Null wertneutral ist:
genau dann, wenn alle Ziffern außer der letzten null sind -- und
das nächste Kapitel macht daraus mit der Horizontableitung einen
Einzeiler.

Wie sich das für eine konkrete Zahl anfühlt, erzählt das
Interludium „Die Lebensgeschichte der Sieben“.
